% !TEX root = ../main.tex
\section{Introduction}
\label{sec:introduction}

Clarity, the smart contract language of Stacks, is interpreted rather than compiled, and a deployed contract is stored on chain as the source text that the node executes \cite{sip002,hiro2024nowhere}. Anyone can therefore read any contract on Stacks mainnet through a block explorer or the Stacks Blockchain API \cite{hiroapi}, and a public repository mirrors the source of every deployed contract on GitHub \cite{boomcrypto}. On Ethereum, by contrast, the source of a contract is visible only when its developer has chosen to publish and verify it. Developers who want to build on an existing Clarity contract have two options. They can call the deployed contract by its principal with \code{contract-call?}, or they can copy its source and deploy a new instance under their own address. Clarinet, the standard development tool, already supports the first option: \code{clarinet requirements add} downloads a mainnet contract into a project so that local tests can call it \cite{clarinet}.

What Stacks lacks is the layer that npm, crates.io and pkg.go.dev provide on top of their code: search, documentation, version lineages, publisher identity, and an ordering that places the contracts other people rely on above copies, tests and spam. Package registries usually order search results with a popularity signal. The Go package site, for example, multiplies the text relevance of a module by the logarithm of the number of packages that import it \cite{pkgsite}. Such counts are cheap to inflate. In April 2024 Sonatype reported about 14,000 npm packages published from several accounts to collect rewards from the tea protocol \cite{sonatype2024tea}, which scores open-source packages by their position in the dependency graph \cite{tea2022}. By November 2025 Amazon Inspector had tied more than 150,000 npm packages to the same campaign; the packages declared dependencies on each other, so that installing one pulled in the rest \cite{aws2025tea}. A Stacks registry whose ranking decides which contract a developer installs will attract the same behavior, and the damage from a bad recommendation is larger, because a referenced contract can hold or move real assets.

Earlier work ranked Ethereum wallets with PageRank. Adaptive Weighted PageRank (AWP) weights ERC-20 transfers by their value and by a logistic decay in their age \cite{do2023awp}. EndorseRank weights each owner--spender edge by the latest ERC-20 allowance between them, on the view that an approval is an explicit grant of spending authority rather than a payment that may have many motives \cite{kwon2026endorserank}. This paper asks whether the same idea can rank Clarity contracts. Two problems stand in the way.

The first problem is that Stacks has no allowance. The SIP-010 fungible token trait defines \code{transfer} but no \code{approve}, \code{allowance} or \code{transferFrom} \cite{sip010}. A user lets a contract move their tokens by calling that contract directly; the contract acts within the user's own transaction, and the user limits what may leave their account with post-conditions, which the protocol checks after execution \cite{sip005}. Since Clarity~4 a contract can also limit its own outflows while it runs code it does not control, through the asset allowances of \code{as-contract?} and \code{restrict-assets?} \cite{sip033}. Together these mechanisms record most of what an allowance records, but each is scoped to a transaction or to a block of code instead of being stored as a standing balance. Section~\ref{sec:mapping} works out which of them can play the role of the allowance edge and how.

The second problem is that PageRank alone does not stop Sybil attacks. Cheng and Friedman proved that no reputation function that treats node identities symmetrically can be Sybil-proof \cite{cheng2005sybilproof}, and link farms raise the PageRank of a target page by surrounding it with pages that link to it \cite{gyongyi2005alliances}. Resistance comes from asymmetry, as in TrustRank and SybilRank, where the walk starts or restarts only at trusted nodes \cite{gyongyi2004trustrank,cao2012sybilrank}. We show in Section~\ref{sec:method} that on a graph in which wallets point to the contracts they authorize, PageRank with a uniform restart gives every wallet the same weight, so the part of a contract's score that comes from wallets becomes a fractional count of the wallets that use it. Moving EndorseRank to Stacks without changing the restart would therefore be about as easy to game as counting callers.

The paper makes four contributions.

\begin{itemize}
\item \textbf{A mapping of the allowance relation to Stacks:} we compare ERC-20 approvals with post-conditions, Clarity~4 asset allowances, legacy \code{as-contract}, extension maps and static references, and we choose post-conditions for wallet-to-contract endorsements and code-level delegation and dependency for contract-to-contract edges. We state which post-conditions count as authorizations. AWP needs no replacement, because Stacks records asset transfers.
\item \textbf{\PCER:} a ranking over contracts only, in which wallets enter through the restart distribution. It keeps the latest-state rule of EndorseRank, weights each wallet by the fees it has paid, and propagates score along delegation and dependency edges. Three propositions describe how it, and EndorseRank moved to Stacks with a uniform restart, respond to Sybil attacks, and a synthetic example illustrates them.
\item \textbf{A registry and command-line client:} the Clarity Registry and \cpm\ separate referencing a deployed contract (\code{cpm add}) from copying its source (\code{cpm create}, \code{cpm fork}), take publisher scopes from deployer addresses and BNS names, treat versions as lineages declared by the publisher, group clones by normalized source hashes and \code{contract-hash?}, and pin code hashes in a lockfile.
\item \textbf{An evaluation protocol fixed in advance:} it compares \PCER\ with raw counts, dependency PageRank, AWP and EndorseRank with a uniform restart on Stacks mainnet, measuring the cost an attacker pays to reach the top of the ranking, agreement with adoption in a later window, and the share of clones among search results. Raw counts are included because in our ERC-20 study they matched held-out approvals better than both PageRank variants \cite{kwon2026endorserank}.
\end{itemize}

Section~\ref{sec:background} describes Stacks, Clarity and the mechanisms that control asset movement, together with the PageRank and registry background. Section~\ref{sec:related} reviews related work. Section~\ref{sec:mapping} maps allowances to Stacks, Section~\ref{sec:method} defines \PCER, and Section~\ref{sec:registry} describes the registry and \cpm. Section~\ref{sec:evaluation} sets out the evaluation, Section~\ref{sec:discussion} discusses limitations, and Section~\ref{sec:conclusion} concludes.
